\documentclass[10pt,a4paper]{amsart}
\usepackage[margin=19mm]{geometry}
\usepackage{amsmath,amssymb,mathtools}
\usepackage[scr=rsfs,cal=euler]{mathalfa}
\usepackage{xfrac}
\usepackage[unicode,hidelinks,bookmarks=false]{hyperref}
\newcommand{\ie}{i.e.\ }
\newcommand{\cf}{cf.\ }
\newcommand{\resp}{respectively\ }
\newcommand{\Subsection}{\subsection}
\providecommand{\nobreakdash}{\nobreak\hbox{-}}
\setlength{\emergencystretch}{2em}
\multlinegap0pt
\usepackage{mathtools}
\usepackage{enumerate}
\usepackage{amssymb}
\usepackage{xcolor}
\newtheorem{theorem}{Theorem}
\newcommand{\R}{\mathbb{R}}
\newcommand{\abs}[1]{\left\lvert#1\right\rvert}
\newcommand{\di}{\,\mathrm{d}}
\newcommand{\fhat}[1]{\widehat{#1}} 
\newcommand{\fourier}{\mathcal{F}}
\newcommand{\norm}[1]{\left\lVert#1\right\rVert}
\newcommand{\smallnorm}[1]{\lVert#1\rVert}
\title{On $L^\infty$ stability for wave propagation \\ and for linear inverse problems}
\author{Rima Alaifari, Giovanni S. Alberti, Tandri Gauksson}
\date{Source: arXiv:2410.11467v2 (CC BY 4.0)}
\begin{document}
\maketitle

\noindent\emph{Source and licence.} On $L^\infty$ stability for wave propagation \\ and for linear inverse problems. Version arXiv:2410.11467v2 (CC BY 4.0), distributed by arXiv under the Creative Commons Attribution 4.0 International licence, http://creativecommons.org/licenses/by/4.0/, as recorded in the arXiv licence metadata for this exact version and shown on the abstract page https://arxiv.org/abs/2410.11467v2. Full licence notice: \texttt{licenses/arxiv-2410.11467-CC-BY-4.0.txt}.\par\smallskip

\noindent\textit{Editorial scope.} Counted unit: the article's theorem thm:regfouriermult\_pwuniformapprox (source0.tex lines 795--802). The article's complete original proof of it is delivered with it (source0.tex lines 803--905, one \texttt{proof} environment of 103 lines). No mathematics is written, repaired or re-derived: the statement and the proof are the article's own lines, in the article's own notation, and no retained line is substituted or re-flowed.  Retained uncounted as the source-local context the proof needs: lines 677--696.  Nothing was omitted from any retained span, so the package carries no omission record.  No retained line contains a citation, so the wrapper carries no bibliography and no bibliography entry.  No retained line contains a figure environment, an image-inclusion command or drawing code, so no image is delivered and the package declares no asset dependency.  Editorial formatting only, in addition: an AMS \texttt{amsart} preamble that restates the article's own declarations and macros used by the retained lines, namely the environments \texttt{theorem} on separate counters, together with the standard packages those lines need.  The retained environments print in this document as Theorem 1 in order of appearance, whereas the article prints the same items with the numbering of its own full document.  The complete native source of the article as distributed in the arXiv e-print for this exact version is bundled unchanged as \texttt{sources/166860/source0.tex}.
\begin{theorem}[Pointwise approximation in \(L^2(\R^d)\)]
\label{thm:regfouriermult_pwl2approx}
    Let \(k,\fhat{h}\in L^1(\R^d)\cap L^2(\R^d)\) be two functions satisfying \(\int k = \int \fhat{h} = 1\).
    Define the families \((k_\alpha)_{\alpha>0}\) and \((h_\beta)_{\beta>0}\)
    through dilation,
    \begin{align}
    \label{eq:approxiddilate_k}
        k_\alpha(x) &= \alpha^{-n} k(x/\alpha)\\
    \label{eq:approxiddilate_h}
        \fhat{h}_\beta(x) &= \beta^{-n} \fhat{h}(x/\beta),\qquad h_\beta = \mathcal{F}^{-1}(\fhat{h}_\beta).
    \end{align}
    Then, for every \(f\in L^2(\R^d)\) we have
    \begin{equation*}
        \norm{T_{\alpha,\beta}f - f}_{L^2(\R^d)}\xrightarrow[\alpha,\beta\to 0]{} 0
    \end{equation*}
    and
    \begin{equation*}
        \norm{B_{\alpha,\beta}f - Bf}_{L^2(\R^d)}\xrightarrow[\alpha,\beta\to 0]{} 0.
    \end{equation*}
\end{theorem}

\begin{theorem}[Pointwise approximation in \(L^\infty(\R^d)\)]
\label{thm:regfouriermult_pwuniformapprox}
    Define the families \((k_\alpha)_{\alpha>0}\) and \((h_\beta)_{\beta>0}\) as in Theorem \ref{thm:regfouriermult_pwl2approx}.
    Then, for any \(f\in L^2(\R^d)\) with \(\fhat{f}\in L^1(\R^d)\) we have
    \begin{equation*}
        \norm{B_{\alpha,\beta}f - Bf}_{L^\infty(\R^d)}\xrightarrow[\alpha,\beta\to 0]{} 0.
    \end{equation*}
\end{theorem}

\begin{proof}
    First note that \(\kappa = \fhat{k}\) is a continuous function by the Riemann-Lebesgue lemma, and that
    \begin{equation*}
        \kappa_\alpha(\xi)
        = \kappa (\alpha\xi)
        \xrightarrow[\alpha\to 0]{}
        \kappa(0)
        = \int k(x)\di x = 1,
    \end{equation*}
    and hence
    \begin{equation*}
        \kappa_\alpha(\xi)\fhat{f}(\xi)
        \xrightarrow[\alpha\to 0]{} \fhat{f}(\xi)
    \end{equation*}
    for almost every \(\xi\in\R^d\).
    Moreover, we have
    \begin{equation*}
        \abs{\kappa_\alpha(\xi)\fhat{f}(\xi)}
        \leq
        \norm{\kappa_\alpha}_{L^\infty(\R^d)}
        \abs{\fhat{f}(\xi)}
        =
        \norm{\kappa}_{L^\infty(\R^d)}
        \abs{\fhat{f}(\xi)}
    \end{equation*}
    for all \(\xi\),
    and since \(\fhat{f}\) is integrable,
    we have
    \begin{equation*}
        \smallnorm{\kappa_\alpha \fhat{f} - \fhat{f}}_{L^1(\R^d)}
        \xrightarrow[\alpha\to 0]{}0
    \end{equation*}
    by the dominated convergence theorem.
    
    Assuming \(\smallnorm{\mu}_{L^\infty(\R^d)}=1\) for simplicity,
    we get:
    \begin{align*}
        \smallnorm{B_{\alpha,\beta}f - Bf}_{L^\infty(\R^d)}
        &=
        \smallnorm{
        \fourier^{-1}
        \left[
        \mu \kappa_\alpha \fourier (h_\beta f)
        -
        \mu \fourier f
        \right]
        }_{L^\infty(\R^d)}\\
        &\leq
        \smallnorm{\mu}_{L^\infty(\R^d)}
        \smallnorm{
        \kappa_\alpha \fourier (h_\beta f)
        -
        \fhat{f}  
        }_{L^1(\R^d)}\\
        &=
        \smallnorm{
        \kappa_\alpha \left( \fhat{h}_\beta *\fhat{f}\right)
        -
        \fhat{f}
        }_{L^1(\R^d)}\\
        &\leq
        \smallnorm{
        \kappa_\alpha \left( \fhat{h}_\beta *\fhat{f}\right)
        -
        \kappa_\alpha \fhat{f}
        }_{L^1(\R^d)}
        +
        \smallnorm{
        \kappa_\alpha \fhat{f}
        -
        \fhat{f}
        }_{L^1(\R^d)}\\
        &\leq
        \smallnorm{\kappa_\alpha}_{L^\infty(\R^d)}
        \smallnorm{
        \fhat{h}_\beta *\fhat{f}
        -
        \fhat{f}
        }_{L^1(\R^d)}
        +
        \smallnorm{
        \kappa_\alpha \fhat{f}
        -
        \fhat{f}
        }_{L^1(\R^d)}\\
        &\leq
        \smallnorm{\kappa}_{L^\infty(\R^d)}
        \smallnorm{
        \fhat{h}_\beta *\fhat{f}
        -
        \fhat{f}
        }_{L^1(\R^d)}
        +
        \smallnorm{
        \kappa_\alpha \fhat{f}
        -
        \fhat{f}
        }_{L^1(\R^d)}
        \xrightarrow[\alpha,\beta\to 0]{}0.
    \end{align*}
    Here we have used the fact that \((\fhat{h}_\beta)_{\beta>0}\) is an approximation to the identity.
    This concludes the proof.
\end{proof}

\end{document}
